This paper demonstrates that county-based congressional representation is both technically feasible and normatively justified as an alternative to traditional redistricting. Using 2020 census data for all 3,142 U.S. counties, we find that 25 counties (70 million people, 21\% of the U.S.) exceed twice the ideal district size and currently span 92 congressional districts. Applying the Huntington-Hill method to allocate seats directly to these qualifying counties achieves representation equality comparable to the current state-based system (CV = 0.100 vs. 0.080) while eliminating fragmentation for one-fifth of the population.

Our central empirical finding is that \textbf{county fragmentation is bipartisan}: 16 qualifying counties are in Democratic-controlled states (46M people), 7 in Republican-controlled states (21M people). Democratic legislatures fragment their own urban strongholds---Los Angeles, Chicago, New York City---for strategic purposes unrelated to partisan advantage: VRA compliance, incumbent protection, suburban vote distribution, and traditional redistricting practices. This demonstrates that fragmentation reflects \textbf{state control over local representation}, not simply Republican gerrymandering of Democratic cities. Both parties prioritize state-level political goals over local community coherence.

We frame county-based representation as fundamentally about \textbf{local autonomy}: natural political communities should not be fragmented for state convenience. Counties are centuries-old governance units with institutional capacity, economic integration, and civic identity. Allowing state legislatures to carve counties into pieces for congressional redistricting violates subsidiarity and creates accountability problems. If a county governs itself locally, it should represent itself nationally.

The partisan reality is inescapable: county-based representation would currently advantage Democrats by an estimated 15-25 seats because large counties are urban and Democratic-leaning. We acknowledge this openly but contextualize it in three ways. First, partisan composition is historically contingent---Orange County and Maricopa County were Republican strongholds a generation ago. Second, current gerrymandering also has partisan bias favoring Republicans; county-based representation levels the playing field for 21\% of the population. Third, principled reforms often have asymmetric impacts---\emph{Reynolds v. Sims} massively advantaged Democrats by empowering urban areas, yet the Court ruled ``one person, one vote'' was constitutionally required regardless of partisan consequences.

Implementation requires innovation. Counties would conduct internal redistricting to divide their allocated seats, maintaining VRA compliance through majority-minority subdistricts while respecting natural county boundaries. House procedures would accommodate county delegations, either through individual representatives (13 from LA County) or weighted voting (fractional votes for small populations). These challenges are significant but not insurmountable if the normative case is accepted.

County-based representation offers unique advantages over alternative reforms. Independent commissions still draw arbitrary boundaries every decade. Algorithmic redistricting creates ``fair'' boundaries with no community basis. Multi-member districts with proportional representation use arbitrary geographic regions. County-based representation combines \textbf{natural boundaries, local autonomy, and proportional fairness} in one system. It eliminates redistricting trauma, gerrymandering opportunities, and the fundamental problem of drawing artificial lines.

The broader implication is that \textbf{redistricting is an unnecessary problem}. We accept decadal boundary disputes as inevitable, but county-based representation reveals an alternative: use fixed, natural boundaries instead of arbitrary, manipulable ones. The principle applies beyond counties---media markets, economic regions, or other natural communities could serve similar functions. What matters is that \emph{boundaries follow communities, not the other way around}.

County-based representation will not be adopted soon. It would require constitutional amendment (Article I, Section 2 prescribes districting), face resistance from state legislatures losing control, and encounter skepticism from those benefiting from current gerrymanders. But as a thought experiment, it clarifies what we sacrifice by accepting redistricting: we fragment natural communities, create accountability problems, invite manipulation, and generate endless litigation---all to maintain arbitrary boundaries redrawn every decade.

If we take seriously the principle that natural communities merit representation, and if we believe local autonomy should extend to federal representation, then county-based representation deserves consideration despite partisan implications. The question is not whether it advantages one party now, but whether it respects natural community boundaries and eliminates manipulation going forward. By that standard, county-based representation succeeds: it preserves the political identity of 70 million Americans while maintaining representational equality and providing a permanent, manipulation-proof alternative to redistricting.

The path forward requires further analysis: demonstrating actual enacted plan fragmentation through spatial overlay, testing VRA compliance through county-internal redistricting models, and assessing public acceptability through surveys. But the core principle is clear: Los Angeles County, Cook County, and Harris County should not be carved into pieces by distant state capitals. They are natural political communities that merit autonomous representation. If that principle is accepted, implementation details can follow.
